(1) $\sigma_y=\begin{pmatrix}0&-i\\i&0\end{pmatrix}$，故 $i\sigma_y=\begin{pmatrix}0&1\\-1&0\end{pmatrix}$。在基向量实数表示下 $\Theta|\uparrow\rangle=-|\downarrow\rangle$、$\Theta|\downarrow\rangle=|\uparrow\rangle$。由于 $\Theta$ 反线性，复共轭作用在第二次运算上也必须计入，最终 $\Theta^2=(i\sigma_y)(i\sigma_y)^*=-I$。
(2) 若 $H|\psi\rangle=E|\psi\rangle$，且 $H$ 时间反演不变，则 $H\Theta|\psi\rangle=\Theta H|\psi\rangle=E\Theta|\psi\rangle$，其中厄米本征值 $E$ 为实。反幺正内积恒等式 $\langle\Theta\phi|\Theta\chi\rangle=\langle\chi|\phi\rangle$，取 $\phi=\psi,\chi=\Theta\psi$，得 $\langle\Theta\psi|\Theta^2\psi\rangle=\langle\Theta\psi|\psi\rangle$；利用 $\Theta^2=-I$ 推出该内积等于其相反数，因此 $\langle\Theta\psi|\psi\rangle=0$，同能伙伴正交。
(3) 在 $\sigma_z$ 本征基，$H_Z=-\mu B\sigma_z$ 的两个能量为 $-\mu B$、$+\mu B$，分裂大小 $\Delta E=2\mu|B|$。时间反演使 $\Theta\sigma_z\Theta^{-1}=-\sigma_z$，固定外磁场 $B$ 不随粒子态时间反演自动翻转，因此 $\Theta H_Z(B)\Theta^{-1}=H_Z(-B)\ne H_Z(B)$，原定理不再保证简并。
(4) $B\to0$ 时自旋Zeeman分裂趋零，若其余项仍时间反演对称，Kramers 成对结构恢复。自旋零常见时间反演满足 $\Theta^2=+I$，前述证明的最后一步不再给出内积必须为零；因此简并需要其他空间、群论或动力学对称性，不能套用自旋半的反幺正论证。
